\section{System Design}

\subsection{Hardware-aware Fine-grained Tile Quantization Scheme}

Existing work~\cite{quantization-hurts-reasoning,quantization-meets-reasoning} has shown that quantization errors significantly degrade model performance in challenging tasks like mathematical reasoning. However, due to stringent on-device resource constraints, full-precision models remain infeasible, making fine-grained quantization essential to maintain accuracy.

Unfortunately, implementing efficient dequantization-based GEMM kernels under fine-grained quantization on mobile NPUs introduces substantial system challenges. We identify two primary issues:

\begin{itemize}
    \item mismatch between the weight layout expected by the matrix unit and conventional group quantization layout;
    \item suboptimal utilization of the wide vector registers caused by small group sizes.
\end{itemize}

To overcome these limitations, we propose a novel tile quantization scheme that incorporates two components:

\begin{itemize}
    \item a tile-based quantization strategy designed to align with the matrix unit's inherent data layout;
    \item a post-quantization weight permutation method that maximizes utilization of the vector unit's processing capabilities.
\end{itemize}

% PART1 Begin (Outdated)

% \ting{too detailed for the beginning. What is research related?}
% In this work, we mainly consider 4-bit quantization of weights, as this is the mainstream configuration for deploying large language models on end devices with relatively limited memory. Following the settings of llama.cpp, we employ group-wise symmetric post-training quantization with a group size of 32, in which 32 quantized elements share the same FP16 scale. The fine-grained quantization group may be less friendly to hardware, but it preserves better precision compared to quantization schemes that use larger quantization groups.

% \todo{improve cohesion(?)} We choose to dequantize the quantized integer weights to FP16 for GEMM mainly for the following reasons: The HMX hardware has support for FP16 matrix multiplication; using FP16 can avoid precision loss caused by outliers without the need to introduce complex QAT pipelines; and it can flexibly support quantized weights of different bit widths.

% To utilize FP16 HMX, specific memory layout transformations must be applied to adapt to the target layout (as shown in Figure~\ref{fig:hmx-fp16-crouton-tile-layout}). For input activation and output, their layout transformations must be performed at runtime. For compatibility reasons, we set the input and output as FP32 type. Figure~\ref{fig:dequantize-based-gemm} demonstrates the mixed-precision GEMM procedure based on dequantization. The input activation blocks are loaded from the global memory into the HVX registers. After the type conversion from FP32 to FP16 and shuffling, we store these blocks in the on-chip memory. The weight blocks, after being dequantized, are also stored in the VTCM in the target layout. The results of FP16 HMX computation are located in the VTCM, followed by type conversions, transforming back to the regular layout, and finally stored in the global DDR memory.

% \begin{figure}
%     \centering
%     \includegraphics[width=0.9\linewidth]{figures/dequantize-based-gemm.png}
%     \caption{Dequantization-based GEMM workflow \ting{I suggest you make graphs first. It should be the first step rather than writing details.}}
%     \label{fig:dequantize-based-gemm}
% \end{figure}

% PART1 End

% \subsubsection{Hardware-aware multi-stage weight reorder}

% Compared to the layout transformation of inputs and outputs, permuting weights’ layout can be performed offline in advance. Since the number of operations required for transforming the weight data during the generation phase of LLM is much larger than that for the input activation and the output, and runtime permutation must be carried out using the relatively weak computation units, it is crucial to prearrange the weights layout offline. (\todo{Our data shows that after introducing the HVX-based layout transformation into the core on-chip HMX computing code, the overall matrix-multiplication throughput has decreased by xxx\%}). We present a multi-level weight rearrangement tailored to the hardware characteristics. The first level is oriented towards the data layout required by the HMX, while the second level is used to adapt the fine-grained quantization groups to the wider HVX registers.

% \paragraph{HMX-oriented static weight reorder}

% The desired weight layout targeting FP16 HMX is illustrated in Figure~\ref{fig:hmx-fp16-gemm-tile-order}. As FP16 HMX performs vector-dot at the tile level, the weight tiles are arranged in column-major order, similar to the regular weight layout (normal linear operations can be expressed as $Y = XW^T$). However, within each tile, the adjacent two rows should be shuffled to satisfy the HMX layout. We carefully consider these requirements and eventually permute the weights data into a nested zigzag order.

% \begin{figure}
%     \centering
%     \includegraphics[width=0.75\linewidth]{figures/hmx-fp16-gemm-tile-order.png}
%     \caption{The memory order of weight matrix for HMX computation. The accumulation functionality of HMX unit favors column-majored tiles, however, elements in each tile are approximately row-majored.}
%     \label{fig:hmx-fp16-gemm-tile-order}
% \end{figure}

\subsubsection{Tile-Group Quantization}

In conventional quantized GEMM, weight matrices are typically stored in column-major layout, which aligns with the vector dot-product operations used in CPU-based matrix multiplication, such as in the llama.cpp CPU backend. The weights are divided into contiguous quantization groups — typically of size 32 — along the column dimension. Within each group, the values are quantized, and the resulting integer weights, along with their corresponding scale and zero-point parameters, are stored interleaved in memory, preserving the original column-wise ordering of the matrix.

However, on NPUs with special matrix units, the conventional group layout is often misaligned with hardware requirements. As illustrated in Figure~\ref{fig:dequant-layout-mismatch}, elements that are contiguous in the conventional layout become scattered in on-chip TCM. For SIMD vector units, such non-sequential access patterns are problematic. Although modern vector engines provide gather/scatter operations to alleviate scattered accesses, these operations remain expensive. Simply transposing the weight matrix does not resolve the mismatch, as the complex multi-level data layout expected by the matrix unit still results in noncontiguous memory access.

\begin{figure}[h]
    \centering
    \includegraphics[width=0.8\linewidth]{figures/dequant-layout-mismatch.pdf}
    \caption{A simplified illustration of the mismatch between the quantization group layout and HMX tile layout.}
    \label{fig:dequant-layout-mismatch}
\end{figure}

To address this, we first permute the weights into the layout expected by the matrix unit, and then apply round-to-nearest quantization group by group. For a group size of 32, this method effectively performs group quantization in units of $2\times16$ tiles. Given that pretrained weights in typical models approximately follow a zero-mean Gaussian distribution, quantizing within these reshaped tile groups does not significantly alter the statistical properties within each group compared to conventional grouping. Therefore, the resulting quantization error remains comparable.

Specifically, we arrange the weights before quantization according to the layout shown in Figure~\ref{fig:hmx-fp16-crouton-tile-layout}, which is hierarchically structured into two levels: an outer column-major ordering of tiles, matching the tile-level inner product operation of the matrix unit, and an inner shuffling of every two rows within each tile. We then quantize the weights group-wise in the new memory order.

\subsubsection{Coalescing Quantization Groups for Wide Vector Accesses}

By default, quantized weights are stored in an Array of Structures (AoS) layout. Taking Q4\_0 symmetric quantization as an example, each group of 32 elements consists of 16 bytes of INT4 quantized values and 2 bytes of FP16 scale values, with quantized values and scales interleaved in memory. Since memory access on the NPU architecture relies heavily on software-managed local 1D or 2D prefetching, we avoid the Structure of Arrays (SoA) layout, where quantized values and scales reside in separate large contiguous arrays, to better align with the hardware's preferred access pattern.

However, fine-grained quantization groups introduce a mismatch with the native vector processing granularity: A single quantization group is too small to fill a 128-byte wide vector register. Accessing such small groups would require multiple memory operations or additional instructions to merge data from multiple registers, resulting in inefficient memory bandwidth usage and computational overhead.

To solve the issue, we coalesce 8 quantization groups into a larger super-group and reorganize its content such that the INT4 values from 256 consecutive elements occupy exactly one full HVX register. This process is illustrated in Figure~\ref{fig:dequant-group-repack}.

% (\hzx{In the literature, this kind of layout is referred to as AoSoA layout.})

\begin{figure}
    \centering
    \includegraphics[width=0.7\linewidth]{figures/dequant-group-repack.pdf}
    \caption{Repacking 8 fined-grained quantization groups into a super-block. The INT4 quantized values fit in a vector register.}
    \label{fig:dequant-group-repack}
\end{figure}

% \subsubsection{Weight dequantization pipeline}

% \begin{figure}
%     \centering
%     \includegraphics[width=1.0\linewidth]{figures/simple-decoding-pipeline.png}
%     \caption{A simple weight load-dequantize pipeline during the decoding phase, where the execution time of the HMX unit is negligible. TODO: further overlap HMX compute during the prefill stage.}
%     \label{fig:simple-decoding-pipeline}
% \end{figure}

\subsection{LUT-Based Computations}

Given the limited general-purpose computing performance of the vector unit, we propose using generalized look-up table (LUT) instructions to replace complex computations, thereby reducing instruction count and computational overhead. LUT-based computation is particularly effective for accelerating key operations in test-time scaling workloads, such as the exponential function in Softmax and the dequantization process.

\subsubsection{Fast Softmax via Vector Gather}

Test-time scaling methods typically increase sampling parallelism, leading to larger batch sizes and longer context lengths. We analyze the impact of these scaling factors on the major operators in transformer-based LLMs during generation:

\begin{itemize}
    \item GEMM. Based on previously described NPU hardware characteristics, moderately increasing batch size in test-time scaling workloads does not substantially increase GEMM latency. Moreover, GEMM latency is independent of context length.
    \item Misc. Ops. For operators such as activation functions, LayerNorm, residual Add, and RoPE, although their computational overhead is roughly proportional to input size, we neglect their impacts due to their small computation and memory access volumes.
    \item Attention. The theoretical computational complexity of Attention scales with both batch size and context length, making it a potential performance bottleneck in test-time scaling scenarios.
\end{itemize}

We implement FlashAttention~\cite{flash-attention2} on the Hexagon NPU using FP16 HMX and measure its latency composition at a prompt length of 4096 under various input batch sizes (query lengths), as shown in Figure~\ref{fig:fa-latency-breakdown}. The results indicate that matrix multiplication contributes little to overall latency, whereas Softmax dominates Attention execution time as the query length increases.

\begin{figure}
    \centering
    \includegraphics[width=0.8\linewidth]{figures/fa-latency-breakdown.pdf}
    \caption{FlashAttention latency breakdown on Hexagon NPU. We use Qwen2.5-1.5B and prompt length is set to 4096.}
    \label{fig:fa-latency-breakdown}
\end{figure}

Our analysis shows that the primary bottleneck of on-chip Softmax lies in the exponential computation, which must be applied to $\Theta(N_q\times N_{kv})$ elements. Adding to the issue, these expensive exponential operations must be executed on the HVX, which lacks dedicated hardware support for special math functions. Following common practice, we replace \texttt{exp} with \texttt{exp2} and absorb the coefficient $\log_2{e}$ in the $QK^T$ scaling factor $\frac{1}{\sqrt{d}}$. For an input element $x$ decomposed into integer part $k$ and fractional part $f$, $2^f$ is approximated using a Taylor series polynomial expansion, while $k$ is directly added to the exponent field of $2^f$'s IEEE-754 representation. However, polynomial evaluation involves sequential dependencies, limiting instruction-level parallelism under the VLIW architecture.

To alleviate the exponential computation bottleneck, we explore replacing explicit exponent calculation with a precomputed lookup table (LUT). The HVX provides the \texttt{vgather} instruction, which can gather values from scattered locations in the TCM into a contiguous 128 byte TCM region. Although \texttt{vgather} can implement large LUTs, using LUTs for $\exp$ remains challenging: storing $2^{32}$ elements for 32-bit floats is impractical. Furthermore, \texttt{vgather} itself introduces substantial latency — 24 to 48 instruction packets on Hexagon V75, so its usage must be minimized.

To enable practical LUT-based $\exp$, we design the following approach. First, we extensively use FP16 throughout FlashAttention, with the on-chip computation process outlined in Algorithm~\ref{alg:fp16-fa-on-chip}. The matrices $S, P, O$ and the vectors $\vec m, \vec l$ are stored in 16-bit floats, with both the input and output of the $\exp$ computation in 16-bit floats. In particular, FP16 HMX uses higher-precision floating-point numbers for accumulation internally, and we upcast elements to 32-bit precision for critical operations such as row-wise summation of matrix $P$.

Using 16-bit inputs and outputs restricts the LUT to 65536 entries, requiring 128 KiB of storage, which fits within the TCM. A variant of \texttt{vgather} supports gathering 64 2-byte elements in one instruction, with a maximum address offset of 65536 bytes. However, 65536 FP16 entries occupy 128 KiB, leaving half of the entries inaccessible with direct addressing. To solve this, we leverage the property of safe softmax~\cite{online-softmax}, which ensures that all inputs to $exp$ are non-positive by subtracting the row-wise maximum $m_i$. Thus, we only store values for $x\le0$, resulting in a LUT with 32768 entries (64 KiB). During LUT-based $\exp$ computation, we ignore the MSB (sign bit) of the FP16 input and left-shift the input by one bit to generate the byte offset required by \texttt{vgather}.

The LUT is precomputed during system initialization, introducing no additional overhead during model inference. It occupies a fixed 64 KiB region in TCM, accounting for only $64KiB/8MiB\approx0.8\%$ of the total TCM capacity, thus minimally impacting TCM availability for other operations.

\begin{algorithm*}[h]
\caption{On-chip computation of ours FP16 FlashAttention (different heads omitted)}
\label{alg:fp16-fa-on-chip}
\KwIn{Head dimension $d$, Number of Query tiles $T_q$, Number of KV tiles $T_{kv}$, Query tile size $B_q$, KV tile size $B_{kv}$}
\KwIn{Matrices $Q_i$ (FP16) $\in \mathbb{R}^{B_q\times d}$, $K_j,V_j$ (FP16) $\in \mathbb{R}^{B_{kv}\times d}$}
Initialize $O_i^{(0)} = (\mathbf{0})\in \mathbb{R}^{B_q\times B_{kv}}$ (FP16), $m = (\mathbf{-\infty})\in\mathbb{R}^{B_q}$ (FP16), $l = (\mathbf{0})\in\mathbb{R}^{B_q}$ (FP16) \;
$S_i^{(j)} = \text{MatMul}(Q_i, K_j^T, \text{AccumType=\textcolor{blue}{FP32}}) \in \mathbb{R}^{B_q\times B_{lv}}$ (FP16) \;
$m_i^{(j)} = \max(m_i^{(j-1)}, \text{rowmax}(S_i^{(j)})) \in \mathbb{R}^{B_q}$ (FP16) \;
$P_i^{(j)} = \text{\textcolor{olive}{LUT\_Exp}}(S_i^{(j)} - m_i^{(j)}) \in \mathbb{R}^{B_q\times B_{kv}}$ (FP16) \;
$l_i^{(j)} = e^{m_i^{(j-1)}-m_i^{(j)}}l_i^{(j-1)} + \text{rowsum}(P_i^{(j)},\text{AccumType=\textcolor{blue}{FP32}}) \in \mathbb{R}^{B_q}$ (FP16) \;
$O_i^{(j)} = \mathrm{diag}(e^{m_i^{(j-1)}-m_i^{(j)}})O_i^{(j-1)}+P_i^{(j)}V_j \in \mathbb{R}^{B_q\times d}$ (FP16, AccumType=\textcolor{blue}{FP32}) \;
\KwOut{$O_i = \mathrm{diag}(l_i^{(T_{kv})})^{-1} O_i^{(T_{kv})}$}
\end{algorithm*}

% \begin{algorithm*}[h]
% \caption{\textcolor{teal}{Implementation of XXX Attention.}}
% \KwIn{Matrices $Q$(FP16), $K$(FP16), $V$(FP16) $\in \mathbb{R}^{N \times d}$, block size $b_q$, $b_{kv}$.}
% \textcolor{olive}{$K = K - \text{mean}(K)$;} \tcp*[r]{Subtracting the mean value across tokens}
% \textcolor{olive}{$(\delta_Q, \hat{Q}) = \psi_Q(Q / \sqrt{d})$, $(\delta_K, \hat{K}) = \psi_K(K)$;} \tcp*[r]{INT8 per - block quant}
% \textcolor{olive}{Divide $\hat{Q}$ into $T_m = N / b_q$ blocks $\{\hat{Q}_i\}$, and divide $\hat{K}$, $V$ into $T_n = N / b_{kv}$ blocks $\{\hat{K}_i\}$ and $\{V_i\}$;}
% \For{\textcolor{purple}{$i$ in $[1, T_m]$}}{
%     \tcp*[r]{Outer loop is paralleled in SMs (stream processors)}
%     \textcolor{olive}{Load $\hat{Q}_i$ and $\delta_Q[i]$ into a SM;}
%     \For{\textcolor{purple}{$j$ in $[1, T_n]$}}{
%         \textcolor{olive}{Load $\hat{K}_j$, $V_j$, and $\delta_K[j]$ into the SM;}
%         \textcolor{olive}{$S^j_i = \text{Matmul}(\hat{Q}_i, \hat{K}_j^T) \times \delta_Q[i] \times \delta_K[j]$;}
%         \textcolor{olive}{$m^j_i = \max(m^{j - 1}_i, \text{rowmax}(S^j_i))$, $\tilde{P}^j_i = \exp(S^j_i - m^j_i)$, $l^j_i = e^{m^{j - 1}_i - m^j_i} + \text{rowsum}(\tilde{P}^j_i)$;}
%         \textcolor{olive}{$O^j_i = \text{diag}(e^{m^{j - 1}_i - m^j_i})^{-1}O^{j - 1}_i + \text{Matmul}(\tilde{P}^j_i\text{.to(FP16)}, V_j, \text{Accum\_type = FP16})$;}
%     }
%     \textcolor{olive}{$O_i = \text{diag}(l^{T_n}_i)O^{T_n}_i$;}
%     \textcolor{olive}{Write $O_i$;}
% }
% \KwOut{$O = \{O_i\}$;}
% \end{algorithm*}


\subsubsection{LUT-Centric Efficient Dequantization}

The runtime HVX dequantization requires careful design to avoid additional overhead. We present an efficient dequantization process based on the HVX lookup table instructions. The \texttt{vlut16} instruction is capable of performing a table lookup in a table of 16 elements for each 8 bit index in a source vector register. Each input byte is transformed into a 16-bit value, therefore \texttt{vlut16} results in a pair of registers.

\paragraph{Fast INT4 to FP16 conversion via table lookup}

Using \texttt{vlut16} instructions, we directly transform 4-bits quantized values into [-8, 7] FP16 values for Q4\_0 quantization scheme, avoiding the conventional mask-unpack-convert instruction sequence. Figure~\ref{fig:lut-based-int4-conversion} demonstrates the comparison of two approaches. For Hexagon NPU prior to V79, all HVX floating-point operations produce results in an internal format called qfloat, which requires extra instructions to convert back to standard IEEE-754 formats. The use of table-lookup eliminates these overheads. This LUT-centric design can easily support different 4-bit encoding schemes (e.g. FP4, NF4~\cite{qlora}, IQ4\_NL used in llama.cpp) simply by adjusting the table contents.

\begin{figure}
    \centering
    \includegraphics[width=0.9\linewidth]{figures/lut-based-int4-conversion.pdf}
    \caption{Converting INT4 quantized values into FP16 numbers via table lookup.}
    \label{fig:lut-based-int4-conversion}
\end{figure}

\paragraph{Scales broadcast via table lookup}

A 128-byte HVX register can accommodate two FP16 quantization groups of size 32. Therefore, the conventional approach is to broadcast scalar scales to the entire vector register and then concatenate two registers for subsequent multiplication with quantized values. However, by using the scales of four groups as LUT contents and applying predefined constant indices, we can achieve the broadcast of four groups of scales with just one \texttt{vlut16} instruction.

% \subsection{Other Potential Algorithmic Optimizations}

% Some of the algorithmic techniques only work for certain TTS methods like step-level beam search, They are not considered as the main contribution.

% \subsubsection{Adaptive Search Path Pruning}

% See "Speculative Rejection" (cite): progressively prunes search paths that are unlikely to be chosen as the final best sequence by the verifier/reward model. This may also work for step-level beam search: terminating extremely long intermediate steps saves computation since those paths are often error prone.

% \subsubsection{Speculative Execution?}

% Specialized for step-level beam search.
